\section{Two hierarchies and their relative spatial origin}\label{sec:line-hierarchy}
\subsection{The imported deterministic construction}
Let $f_n$ be exactly the radius-one hardwired rule of
\cite[Theorem~C and Section~3.5]{DMline}, with its unique empty symbol.
Its literal is the true initial configuration of $U(n)$, doubled at
successive promotions. A promotion executes one machine step.
No external parameter reader is used. Geometry, decoding, interpreter
and native transition are unchanged. Section~3 of that reference
makes the finite table computable from $n$, at fixed $P,T,w$.
A logical step below means one native update.

We use the following numbered statements of \cite{DMline}.
Definition~4.1 defines $i$-sparsity by deleting one translate of
$3w(P^{i-1},T^{i-1})$ in every $w(P^i,T^i)$ window;
$0$-sparsity means emptiness. Lemma~4.2 gives disjoint child witnesses.
Definition~5.1 supplies nominal and served tile observations, whose
spatial support is inside the tile and whose lookahead is at most
$T^j-1$. Proposition~5.3 gives organization for every translated grid,
every history with sparse deviations, true read-only operands and
ordinary unhalted lower ranks; in particular nominally live tiles
with distinct geographical origins cannot overlap.
Lemma~5.24, with six consecutive served positive inputs of one tribe
and two prescribed positive central transitions, gives service of
both central homes for two periods, the correct lower clock and
address residues, and the two served prescribed upper successors.
It assumes a fully protected $(j-1)$-sparse rank-$j$ test.
Lemmas~6.2 and~6.3 give the halted all-start cut and ascent from any
live physical site to a failed centered test after a common burn-in.
Lemma~6.4 gives the clean initialized negative-time hierarchy in the
hardwired family. Lemma~6.5 constructs its deterministic nested tests.
The clock witness in Theorem~C alone would not imply a count of
stationary row laws; we need the spatial-residue conclusion of
Lemma~5.24 and the decoder below.

\subsection{A quantitative mark estimate}
\begin{lemma}[Sparse tests and fixed-rank paths]\label{lem:line-sparse}
Suppose deviations of any native history are contained in iid marks
of rate at most $e$. For every prescribed centered test with half-sides
$M_*(P^{i+1},T^{i+1})$, failure of $i$-sparsity has probability at most
\[
 a_i=A_*V_*^{-2\cdot2^i},\qquad r_i=R_*F_*^i
\]
and is supported in a max-norm ball of radius $r_i$.
Moreover
\[
 \sum_i a_i<\tfrac14,\qquad
 \sum_i(i+1)a_i<\infty,\qquad
 \sum_i10^6r_i^4\sqrt{a_i}<\tfrac12.
\]
For each fixed $K$, a nearest- or diagonal-neighbor path of vertices
with a failed centered test of some rank $0\le i\le K$, starting at
a prescribed point and reaching max-distance $h$, has probability
at most $C_K\exp(-c_Kh)$.
\end{lemma}
\begin{proof}
All tests needed here fit the stated half-sides. One-rank work is in
a $w$-window and its collar has $\lfloor w/5\rfloor$ cells. The nested
tests of \cite[Lemma~6.5]{DMline} have fewer than $w$ cells at their
own scale. Cap ancestry at scale $i+1$ costs at most
$(4T/P+4,6)$ cells before its work collar and rounding
\cite[Lemma~6.3]{DMline}. Covering every translate of a tile through
one point costs at most two further cells per side. All these finite
multipliers are below $M_*$.

Let $b_i$ bound failure on a full doubled box with side vector
$2w(P^i,T^i)$. Its base bound is $b_0\le4w^2e$.
The disjoint witnesses of \cite[Lemma~4.2]{DMline} have at most
$(2P+3)(2T+3)$ possible child positions, hence at most $D_*$ pairs.
Independence on disjoint child boxes gives $b_{i+1}\le D_*b_i^2$.
Induction yields
\[
 b_i\le D_*^{-1}(4D_*w^2e)^{2^i}
       =D_*^{-1}V_*^{-2\cdot2^i}.
\]
Cover a larger prescribed box by a mesh of spacing $w(P^i,T^i)$,
including one mesh step past each boundary. Every moving $w$-window
has its intersection covered by one doubled box. There are at most
$A_*$ such boxes. If all pass, restriction gives the deletion required
in every moving window. At rank zero the same cover bounds every
marked site. The union bound proves the asserted $a_i$; the enclosing
box has radius at most $r_i$. The estimate applies to arbitrary error
values, since sparsity is inherited by subsets of the marks.

Using $2^i\ge i+1$, the first two sums are bounded respectively by
\[
 \frac{A_*}{V_*^2(1-V_*^{-2})},\qquad
 \frac{A_*}{V_*^2(1-V_*^{-2})^2}.
\]
The third is at most
\[
 \frac{10^6R_*^4\sqrt{A_*}}{V_*(1-F_*^4/V_*)}<\tfrac12.
\]
The definitions in \eqref{eq:line-rate} give all three inequalities.

To prove the path assertion, cover a finite such path by its failing
test balls. Greedily retain disjoint balls in decreasing radius.
Every removed ball is contained in a retained triple ball. A simple
chain of adjacent retained triple balls connects the endpoints.
The sum of its radii is at least $h/8$: the diameters are at most $6r$
and a unit jump costs at most one further radius.
The first center has at most $(6r+1)^2\le49r^2$ choices.
Successive centers of radii $r,s$ have at most
$(6r+6s+3)^2\le225r^2s^2$ choices. For an ordered type list the total
placement count is therefore at most the product of $10^6r_i^4$.
Disjoint retained supports make the failure probability at most the
product of their $a_i$. Reserve half of each exponent. Since the
radius sum is at least $h/8$, this half contributes at most
\[
 \exp\left[-\frac h8
       \min_{0\le i\le K}\frac{-\log a_i}{2r_i}\right].
\]
Sum the remaining half over type lists and placements. The resulting
geometric series converges because $\sum_i10^6r_i^4\sqrt{a_i}<1/2$.
This proves the claimed exponential estimate. Only $K$ enters its
exponent. The entire event is a function of the common dominating
marks, independently of which histories might realize discrepancies.
\end{proof}

\subsection{Forgetting on halting}
\begin{lemma}[All-input common-map forgetting]\label{lem:line-forget}
If $U(n)$ halts, a native history whose deviations are contained in
the marks of Lemma~\ref{lem:line-sparse} forgets all starting rows on a
common event with one-site failure probability $C\exp(-ct)$, provided
that with innovations fixed every update, even a marked update, is a
deterministic function of its three old symbols. The constants may
depend on $n$. The bound is for the event that the entire graphical
map at the observed site is nonconstant.
\end{lemma}
\begin{proof}
Let $K$ be the first halted rank. Its true literal forces the protected
strong issue to be empty on every triple, including forged input and
work symbols; all lower ranks are ordinary. For $K\ge1$,
\cite[Lemmas~6.2--6.3]{DMline} therefore apply literally. Their premises
use the native rule off deviations, the unchanged geometry and
true hardwired operands, all of which hold here. Choose the common
burn-in greater than twice the largest radius of the finitely many
centered tests below $K$, plus two. After this burn-in any live vertex
has a failed centered sparsity test below $K$. This inclusion is
simultaneous for every starting row on the same marks.

At a discrepancy of two trajectories at least one symbol is live,
since the empty symbol is unique. If the three parents were identical,
the common deterministic update with fixed innovations would give
identical children. Trace a discrepant parent backwards to half the
elapsed time. For $t$ beyond twice the burn-in plus four, this gives a
nearest- or diagonal-neighbor path of failed tests with time span at
least $\lfloor t/2\rfloor-1$. Lemma~\ref{lem:line-sparse} makes its
probability exponentially small. The same noise-only path is necessary
if \emph{some} pair of starting rows disagrees, so no union over rows
is needed. Finite exceptional times enter the prefactor.

At $K=0$ an unmarked output is always empty. A discrepancy must follow
a backward path all of whose vertices are marked. There are at most
$3^t$ length-$t$ paths and each has probability at most $e^t$, giving
$(3e)^t$. For exact classical replacement, the rank-zero constant rule
in fact coalesces in one update. A marked encoded update need not erase
a discrepancy, so the path bound is used for that implementation.

If an ascent test reads beyond the observation date, independently
extend the innovations to those dates. They can witness the necessary
path without changing the nonconstancy event, which depends only on
the actual backward cone. Finally a union bound over an observed finite
set gives the bound for its whole common map.
\end{proof}

\subsection{Decoding a spatial origin}
\begin{lemma}[Stopped descent and finite decoders]\label{lem:line-decoder}
If $U(n)$ never halts, initialize its clean aligned hierarchy from
\cite[Lemma~6.4]{DMline}. For every $j\ge1$, a finite space-time decoder
at $(x,t)$, $t\ge T^j$, returns the left edge of its true rank-$j$ tile
modulo $P^j$, except with probability at most
\[
 d_j=\sum_{i\ge j-1}a_i,\qquad \sum_{j\ge1}d_j<\infty.
\]
The bound is uniform in $x,t$ and integer spatial shifts of the
initialization. At rank zero, with probability at least
$1-\sum_i a_i>3/4$, the physical root is live with both initialized
structural residues, in particular its spatial address modulo $P$.
\end{lemma}
\begin{proof}
Every rank is ordinary on a nonhalting input. Join the fully occupied
clean negative-time history to any forward history obeying the native
rule off the marks. For $j\ge1$, examine all integer tiles of width
$P^j$ and height $T^j$ containing $(x,t)$, using half-open tiles.
Compute their finite nominal observations from
\cite[Definition~5.1]{DMline}. Return the common left edge modulo
$P^j$ if at least one is nominally live and all nominally live
candidates have that edge; otherwise fail. The decoder is a cylinder
function, possibly reading dates on either side of $(x,t)$.

Start the geometrical ascent of tests with the entire true tile.
Its dimensions are less than $P^{j+1}$ and $T^{j+1}/2$, as required
by \cite[Lemma~6.5]{DMline}. For completeness, if a smaller time span
is $[a,b)$ of length less than $T^k/2$, put
$h=\lfloor a/T^k-1/2\rfloor$; it lies in
$[(h+1/2)T^k,(h+2)T^k)$. If its spatial span is $[c,e')$ of length
less than $P^k$, choose $k'=\lfloor c/P^k\rfloor$ and the two homes
$k',k'+1$. Add the six bottom reads, their work, covering homes and
the prescribed collar. The resulting box again fits the next rank,
by the strict inequalities of that lemma. Every choice is determined
by the true grid and the observation, independently of the noise.

At a sufficiently large finite top rank $r$ choose first test date
$-T^r$. Enlarge $r$ so $4(t+2)<T^r$ and the bottom tile and collar fit.
All six top reads, including their lookahead at most $T^r-1$, are
strictly before zero. The initialized past supplies their common
tribe and both prescribed positive successors. Impose
$(k-1)$-sparsity for these tests at $k=j+1,\ldots,r$.
Apply \cite[Lemma~5.24]{DMline} at the top. It serves both central
homes throughout the covering two periods with their correct child
residues. Its two upper successors supply the actual positive
transitions needed at the next descent. Repeating the same argument,
with those actual served inputs and successors at each rank, and
stopping at $j$, gives a served live observation on the true tile.
A served live observation is nominally live. No recovery of arbitrary
mutable root work data is asserted.

For uniqueness, put the true tile, all overlapping candidates and
their collars in one prescribed test of scale $j$. The size bound was
checked in Lemma~\ref{lem:line-sparse}. Its $(j-1)$-sparsity implies
$j$-sparsity. The classification of \cite[Proposition~5.3(i)]{DMline}
excludes nominally live overlapping candidates with different spatial
origins. In that proposition geographical origin is the full space-time
offset modulo $(P^j,T^j)$; its shifted-child-grid argument uses both
offsets. The decoder's common-edge convention also works using only
the spatial-origin consequence. The true tile supplies existence,
and all live candidates therefore have its spatial residue.
The failed-test union has probability at most
$\sum_{i\ge j-1}a_i$. Negative-time deviations are empty and may be
dominated by an independent extension of the marks. This proves the
uniform bound, and the weighted sum in Lemma~\ref{lem:line-sparse}
proves summability. Continuing the descent through rank one instead
of stopping gives the physical root and both residues by the explicit
conclusion of \cite[Lemma~5.24]{DMline}, with error at most $\sum_i a_i$.
\end{proof}

\subsection{Stationary laws with distinct relative origins}
\begin{proposition}[Continuum many ordered extreme states]\label{prop:line-phases}
Use two autonomous copies of the nonhalting hierarchy on the same
line. Their noises may be dependent, provided each satisfies
Lemmas~\ref{lem:line-sparse}--\ref{lem:line-decoder}.
More generally there may be a fixed positive number of logical
steps per observed round. The joint controller has continuum many
extreme TI stationary states with a single-row Bragg atom at $\pi$.
\end{proposition}
\begin{proof}
Let $G=\varprojlim_j\Z/P^j\Z$. It has cardinality $\frakc$: each digit
may independently be zero or one, while it is a countable product
of finite sets. Fix $d\in G$. For $h\ge1$, let $\hat d_h\in[0,P^h)$ be
its representative. Initialize the first copy with spatial origin
zero and the second with origin $\hat d_h$, on the same time grid.
Neither rule nor noise reads $d$. In an encoded implementation start
at the boundary preceding fanout with encoded centers and default
auxiliaries. Run the actual dynamics. Average its trajectory over
spatial shifts $0,\ldots,h-1$ and time shifts by $0,\ldots,h-1$
complete rounds, extending negative dates arbitrarily for compactness.
Take a common weakly convergent subsequence on every finite window.

Boundary defects in either average are at most $2/h$, so the limit
is TI and invariant under complete-round time translation. The
conditional transition identities against any finite past cylinder
are the true kernels: each initialized experiment uses fresh
innovations, complete-round shifts retain the same schedule phase,
and the proportion of shifts whose cylinder reaches before the
initialization tends to zero. These identities pass to the limit
on cylinders and then by approximation. Thus the limit is the
canonical two-sided trajectory of its stationary boundary row,
with the intermediate kernels inserted in their specified order.
For any stationary row law this trajectory map is affine on finite
cylinders and Borel; use a past boundary and stationarity for
two-sided cylinders. In the quantum implementation all boundary
controller states are diagonal, as proved in
Lemma~\ref{lem:line-physical}, so the same construction applies to the
full basis-valued register and its local logical decode.

Apply the two rank-$j$ decoders at the space-time origin and subtract
the first returned residue from the second modulo $P^j$, failing if
either fails. Common spatial shifts cancel; time shifts do not move
a spatial origin. For $h\ge j$ the correct difference is $d\bmod P^j$.
At fixed $j$ the proportion of starts too early is $O(T^j/h)$, and
on the others the error is at most $2d_j$, with $d_j$ from
Lemma~\ref{lem:line-decoder}. Each decoder is a finite
cylinder of the intermediate history, so the same bound holds in
the limit. Borel--Cantelli gives probability one to the Borel event
$B_d$ that every sufficiently large output is the residue of $d$.
For distinct $d$ these events are disjoint, since a differing residue
remains different at all higher ranks.

Let $f$ on the first lane be $(-1)^a$ on a live symbol with spatial
address $a$, and zero on the empty symbol; for encoded rows use the
locally decoded observable. Since $P$ is even,
Lemma~\ref{lem:line-decoder} gives in the origin-zero experiment,
for every site and date, including round boundaries,
\[
 \mathbb E[(-1)^x f(x)]\ge1-2\sum_i a_i>\tfrac12.
\]
A spatial shift changes the prescribed sign by one common parity.
For every interval length $l$, Jensen's inequality before the
shift mixture gives
\[
 \mathbb E\left|\frac1l\sum_{x=0}^{l-1}(-1)^xf(x)\right|^2
 \ge\tfrac14.
\]
This is a finite row observable. The bound passes through both
averages and the weak limit. The phase-twisted mean ergodic argument
in Section~\ref{sec:framework} shows that its limit is the Bragg
second moment at $\pi$.

The TI stationary controller set is compact convex metrizable in
the locally convex space of its countable real cylinder coordinates.
The metrizable representation theorem
\cite[Theorem~4.2, pp.~316--317]{bishop-deleeuw} therefore represents
our row state by a probability measure carried by its extreme points.
The mixture identity for their canonical trajectories holds on
cylinders by linearity of kernels and stationarity, hence on all
Borel sets by the monotone class theorem. Since $B_d$ has probability
one, almost every representing component also gives it probability
one. The finite Bragg second moments are bounded affine functions
of a row state and converge in every TI state. Bounded convergence
shows that the integral of the limiting moments is at least $1/4$.
A set of positive representing measure has a positive moment, so
some extreme component has both label $d$ and a peak at $\pi$.
A stationary row state has a unique canonical trajectory and cannot
carry two disjoint probability-one labels. Varying $d$ gives
$\frakc$ distinct such extremes; countable coordinates give the
matching upper bound. No uniqueness of the representation, spatial
ergodicity of its components, or conditioning of the dynamics on a
noncausal event is used.
\end{proof}
